\subsection*{Chapter 2: PyTorch and the Machine Underneath}

\paragraph{Exercise 2.1}
Contiguous $(2,3,4)$ has strides $(12,4,1)$. After \code{transpose(0, 2)} the shape is $(4,3,2)$ and the strides
are $(1,4,12)$, over the same storage. \code{.view(-1)} requires that the flattened order equal the storage order,
which no longer holds, so it raises an error. \code{.reshape(-1)} falls back to a copy: one kernel that reads and
writes all 24 elements (in general, $2\times$ the tensor's bytes of traffic plus a new allocation).

\paragraph{Exercise 2.2}
bf16 has 7 explicit mantissa bits, so the spacing of representable numbers in $[1,2)$ is $2^{-7}=0.0078125$. With
round-to-nearest-even, $1+2^{-8}$ rounds back to 1, so any $\delta$ slightly larger than $2^{-8}\approx0.0039$ is needed to
change the value. An update of $10^{-4}$ is lost entirely: the weight never moves. fp32 master weights have
spacing $2^{-23}\approx1.2\times10^{-7}$ near 1, so the update is applied exactly in fp32 and the bf16 copy is
re-derived from the master after each step; small updates accumulate until they cross a bf16 rounding
boundary.

\paragraph{Exercise 2.3}
With $a=\sigma(z)$ and $\mathcal{L}=-[y\log a+(1-y)\log(1-a)]$:
$\partial\mathcal{L}/\partial a = -y/a + (1-y)/(1-a) = (a-y)/(a(1-a))$ and $\partial a/\partial z=a(1-a)$, so
$\partial\mathcal{L}/\partial z=a-y$. For a batch with mean loss, $\nabla_{\vect{w}}=\frac1N\mat{X}^\top(\vect{a}-\vect{y})$
and $\partial/\partial b=\frac1N\sum_i(a_i-y_i)$. This is exactly \code{logistic\_regression\_grads} in the lab
solution, which matches autograd to $10^{-6}$.

\paragraph{Exercise 2.4}
Parameters: $(784\cdot512+512)+(512\cdot256+256)+(256\cdot10+10)=401{,}920+131{,}328+2{,}570=535{,}818$. In fp32:
weights 2.14\,MB, gradients 2.14\,MB, Adam's two moments 4.29\,MB, for a total of 8.57\,MB (16 bytes/parameter). Forward
FLOPs for a batch of 128: $2\times128\times(784\cdot512+512\cdot256+256\cdot10)\approx1.37\times10^8$ (biases and ReLU add
well under 1\%).

\paragraph{Exercise 2.5}
A100: $I^\star=312/2.04\approx153$\,FLOP/byte; a bf16 square GEMM has $I=n/3$, so it becomes compute-bound at
$n\approx460$. RTX~4090: $I^\star=165/1.008\approx164$, so $n\approx490$. The decode critical batch is
$B^\star\approx bI^\star/2=I^\star$ in bf16, so the 4090 needs a slightly larger batch (about 164 versus 153). Both need
far smaller batches than an H100 ($\approx$295), whose compute has outgrown its bandwidth.

\paragraph{Exercise 2.6}
$d=4096$, hidden $4d=16{,}384$, bf16. \emph{$BT=1$}: each GEMM does $2\times4096\times16{,}384\approx1.34\times10^8$ FLOPs and
moves about $2\times4096\times16{,}384\approx134$\,MB (the weight dominates), so $I\approx1$: memory-bound, about 40\,\textmu s each on an
H100. GELU touches 16{,}384 elements, which is negligible but still costs a launch. \emph{$BT=4096$}: GEMM\,1 does $5.5\times10^{11}$
FLOPs and moves $2(4096^2+2\cdot4096\cdot16{,}384)\approx302$\,MB, so $I\approx1{,}800$: compute-bound, $\approx0.56$\,ms at peak; GEMM\,2
is the same. GELU reads and writes $4096\times16{,}384$ bf16 values, 268\,MB, $\approx0.08$\,ms, memory-bound. Fusing GELU into
GEMM\,1's epilogue removes that round trip (the GEMM output is activated in registers before being stored),
saving $\approx0.08$\,ms of $\approx1.2$\,ms, about 7\%. Fusion matters far more for chains of elementwise ops without a GEMM to
hide behind.

\paragraph{Exercise 2.7}
Gradients are 14\,GB in bf16. A ring all-reduce sends and receives $2(P-1)/P\times14\approx24.5$\,GB per GPU for $P=8$. Over
NVLink at 300\,GB/s this takes $\approx82$\,ms, easily hidden behind a 1.5\,s step. Over 100\,Gb/s Ethernet (12.5\,GB/s) it takes
$\approx2.0$\,s, longer than the compute, so communication cannot be hidden and the step at least doubles. The model does
not fit for DDP anyway: 7B parameters $\times$ 16 bytes $=112$\,GB of weights, gradients and Adam state exceeds 80\,GB before
activations. This requires ZeRO/FSDP sharding (\cref{ch:distributed}).

\paragraph{Exercise 2.8}
For 1{,}024 elements, both eager and compiled versions are dominated by per-op overhead: eager time is about
$100\times$ the per-launch cost (hundreds of \textmu s to milliseconds); this is regime (b) of
\cref{fig:pytorch:async}. For $10^8$ elements, each op is a memory-bound kernel taking $\approx0.24$\,ms in fp32 on an
H100, so eager time is about 24\,ms and overhead is hidden; this is regime (a). \code{torch.compile} fuses the chain into one
kernel that reads and writes once. The large tensor speeds up by up to $\approx100\times$ in traffic terms, and the small
tensor speeds up because 100 launches became one; this is regime (c). With \code{reduce-overhead}, CUDA graphs remove
the remaining launch cost.

\paragraph{Exercise 2.9}
At batch 8, the profile is dominated by many small kernels: \code{addmm}/\code{mm} at low intensity (memory-bound,
weight reads dominate), plus optimiser elementwise kernels (\code{add\_}, \code{mul\_}, foreach ops) that touch every
parameter and are memory-bound. At batch 4096, \code{addmm}/\code{mm} in forward and backward take the large majority of
device time and are compute-bound; elementwise ops shrink to a small fraction. Optimiser cost is independent of batch
size, so its share falls as the batch grows.

\paragraph{Exercise 2.10}
Achieved model FLOP/s $=2.1\times10^6\times6\times3\times10^9=3.78\times10^{16}$; peak $=64\times989\times10^{12}=6.33\times10^{16}$; MFU
$\approx60\%$. That is at the very top of what well-tuned dense training achieves, so ask: are padding tokens counted as
tokens? Is $N$ the total or non-embedding count? Are attention FLOPs or recomputation included (that would make it
hardware FLOPs utilisation, HFU, not MFU)? Is the peak taken as the bf16 or the fp8 rate? Is the throughput a
steady-state peak or an average including evaluation and checkpoint stalls?
